\subsection{Estimating the overall reference fraction}
\label{sec:supp-p-estimation}

The parameterization
\[
p_M=E\{\pi_M(W)\},
\qquad
e_M(W)=\pi_M(W)/p_M
\]
is used to state the triangular rate and covariance limit. The adaptive coefficient does not require \(p_M\), because multiplying \(G_{M,h}\) and \(H_{M,h}\) by the same positive scalar leaves their ratio unchanged. For process scaling under an estimated verification law, \(p_M\) can be estimated directly from the observed verification indicators.

\begin{proposition}[Reference-fraction estimation]
\label{prop:p-estimation}
Let
\[
\widehat p_M
=
\frac1M\sum_{i=1}^MR_i.
\]
Under Conditions~\ref{cond:sampling} and \ref{cond:kernel},
\begin{equation}
\frac{\widehat p_M}{p_M}-1
=
O_{\mathbb P}\{(Mp_M)^{-1/2}\}
=o_{\mathbb P}(1).
\label{eq:p-relative-rate}
\end{equation}
Consequently,
\[
\sqrt{M\widehat p_Mh^d}
=
\sqrt{Mp_Mh^d}\{1+o_{\mathbb P}(1)\},
\]
and replacing \(p_M\) by \(\widehat p_M\) in process normalization, multiplier normalization, and the reported local-information rate index leaves the limiting laws unchanged.

If
\[
\widehat p_M^{\pi}
=
\frac1M\sum_{i=1}^M\widehat\pi(W_i)
\]
is formed from cross-fitted propensity predictions and
\[
\frac1M\sum_{i=1}^M
\frac{|\widehat\pi(W_i)-\pi_M(W_i)|}{p_M}
=o_{\mathbb P}(1),
\]
then \(\widehat p_M^{\pi}/p_M\to_{\mathbb P}1\) and the same conclusion holds.
\end{proposition}

\begin{proof}
Since \(E(R_i)=p_M\),
\[
\operatorname{Var}(\widehat p_M)
=
\frac1M\operatorname{Var}(R_i)
\le
\frac{p_M}{M}.
\]
Chebyshev's inequality gives \eqref{eq:p-relative-rate}. The condition \(Mp_Mh^d\to\infty\), together with \(h^d\le1\) eventually, implies \(Mp_M\to\infty\). The square-root expansion follows from differentiability at one. For \(\widehat p_M^{\pi}\),
\[
\left|
\frac{\widehat p_M^{\pi}}{p_M}-1
\right|
\le
\left|
\frac{M^{-1}\sum_i\pi_M(W_i)}{p_M}-1
\right|
+
\frac1M\sum_i
\frac{|\widehat\pi(W_i)-\pi_M(W_i)|}{p_M}.
\]
The first term is \(O_{\mathbb P}(M^{-1/2})\) because \(e_M=\pi_M/p_M\) is uniformly bounded; the second tends to zero by assumption.
\end{proof}
